\documentclass[11pt]{article}
\usepackage[letterpaper,margin=1in]{geometry}
\usepackage[T1]{fontenc}
\usepackage[utf8]{inputenc}
\usepackage{microtype}
\usepackage{xcolor}
\usepackage[colorlinks=true,allcolors=blue]{hyperref}
\usepackage{parskip}

\definecolor{commentgray}{gray}{0.25}
\newcommand{\Status}[1]{\textbf{Status: #1.}}
\newenvironment{reviewcomment}{\begin{quote}\color{commentgray}\itshape}{\end{quote}}

\title{\textbf{Response to the Editor and Reviewers}\\
\large RankCloak: Hiding Synthetic Cryptographic Artifacts in Plain English with Language Models}
\author{Alexander V. Mantzaris\\Submission ID: 0170c7b2-1065-4344-a5d5-b839b51ac22c}
\date{16 August 2026}

\begin{document}
\maketitle

We thank the editor and reviewers for identifying the limitations of the original submission. The revision replaces the partial-pilot narrative with the completed, manifest-bound study. It expands the corpus to 480 public deterministic cryptographic test vectors, three model families, 18 English templates in six categories, and secondary Spanish and Simplified Chinese tests; adds grouped inference, controlled perturbations, automated readability diagnostics, ablations, prompt-topic contrasts, capacity theory, computational overhead, and two neural detector architectures; and narrows the claims in light of severe visible-text fragility and high neural detectability.

The principal positive result is now explicitly conditional: 6,480/6,480 primary serialized payloads recovered under saved-token exact-copy replay. The central adverse results are equally prominent: detokenized-text replay recovered 88/144, common transformations reduced recovery, and matched detector ROC-AUC was 0.990190 for TextCNN-equivalent and 0.999834 for DeBERTa-v3-base. No participant study or deployment experiment was performed. The response below distinguishes completed work, partial responses, and external actions rather than treating every request as fully satisfied.

\section*{Editor}

\subsection*{Accuracy, scale, methodology, and moderated conclusions}
\begin{reviewcomment}
The editor requested substantial experimental and methodological improvement: a larger corpus, statistical analysis, multiple LLM architectures and prompt categories, realistic cryptographic payloads, human readability assessment, stronger neural steganalysis, ablations, failure analysis, overhead measurement, and closer examination of exact-copy replay. Results should be accurate, conclusions moderated, and limitations fully explained.
\end{reviewcomment}

\Status{Substantially addressed}

The revised manuscript incorporates the completed evidence for every computational item listed. Main Table 1 summarizes the 480-artifact design, three models, 18 templates, 6,480 primary payload trials, grouped sensitivities, two secondary languages, and the complete 38,504-unit ledger. Four main figures address capacity/tail cost, transmission fragility, automated readability, and neural detectability; Main Table 2 summarizes ablations, topic effects, and overhead. The Supplement provides full methods, work-unit reconciliation, transformation coverage, all 60 canonical ablation contrasts, all 15 topic contrasts, detector leakage/sensitivity results, theory, timings, and worked failures.

The conclusions no longer present RankCloak as practical, robust, secure, or undetectable. The Discussion treats high detector performance as adverse evidence and concludes that the method is a reproducible measurement framework, not a demonstrated covert channel. The human-evaluation component remains only partially addressed: established automated metrics were added, but there were zero participants and ratings. Deployment was not performed.

\textbf{Locations:} Abstract, p.~1; Main manuscript Results, pp.~5--11, lines 220--390; Discussion, pp.~11--12, lines 392--439; Responsible use and limitations, p.~12, lines 440--461; Supplementary Notes S1--S10, pp.~1--19, lines 5--220.

\subsection*{DOI-assigning code archive}
\begin{reviewcomment}
The editor states that custom code or a new algorithm/pipeline must be deposited in a recognized DOI-assigning repository and linked from Methods or Code Availability.
\end{reviewcomment}

\Status{External action pending}

The revision identifies the public repository and exact Git commit and describes the manifest-bound retained evidence. A DOI has not been supplied or invented during this drafting stage. If the journal requires archival deposition for resubmission, the author must create the DOI-bearing archive and insert the verified DOI before submission.

\textbf{Location:} Data and code availability, p.~12, lines 463--468.

\section*{Reviewer 1}

\subsection*{1. Exact-copy replay and formatting disruption}
\begin{reviewcomment}
The exact-copy replay assumption is highly restrictive. Please add theoretical discussion or experiments addressing minor formatting disruptions during transmission.
\end{reviewcomment}

\Status{Substantially addressed}

We now distinguish saved token IDs from ordinary visible-text transmission. Saved IDs are the strongest replay record, while detokenization/retokenization recovered only 88/144. Raw unmodified text recovered 144/144, but NFKC, quote conversion, whitespace changes, line-ending conversion, character/token edits, markdown copy/paste, paraphrase, and cross-model decoding were damaging or fatal. The final-10\%-tail truncation result is explained narrowly: only declared non-payload tokens were removed. Limited canonicalization did not recover the available transformed cells. The Supplement inventories five directly tested, three partially represented, and three untested requested perturbation classes and avoids causal claims from first-divergence labels.

\textbf{Locations:} Methods, pp.~2--3, lines 80--122; Results and Fig.~2, pp.~7--8, lines 254--279; Discussion, pp.~11--12, lines 398--404; Supplementary Note S4, pp.~6--8, lines 80--87.

\subsection*{2. Evaluation-corpus size}
\begin{reviewcomment}
Two examples per payload class are insufficient; expand the dataset to a substantially larger and more diverse sample.
\end{reviewcomment}

\Status{Fully addressed}

The completed corpus contains 480 artifacts: eight classes with 60 public deterministic artifacts per class. The primary design contains 14,400 work units, including 6,480 payload-bearing trials. Trial-, payload-, and model--payload-group endpoints and Wilson intervals are reported separately so repeated prompt/protocol observations are not mistaken for independently generated artifacts.

\textbf{Locations:} Methods, p.~4, lines 138--160; Results and Table~1, pp.~5--6, lines 220--236; Supplementary Notes S2--S3, pp.~3--6, lines 54--78.

\subsection*{3. Human-perceived naturalness}
\begin{reviewcomment}
Log-probability proxies do not establish human-perceived naturalness; incorporate a participant study or established readability metrics.
\end{reviewcomment}

\Status{Partially addressed}

We added established automated surface diagnostics on 504 balanced computational stimuli, including Flesch Reading Ease, surface flags, prompt similarity, prompt-template bootstrap intervals, and source/held-out model-based contrasts. The manuscript repeatedly states that these are not human ratings. No ethics-authorized recruitment or survey deployment occurred, and there were zero participants and zero ratings. Prepared materials are treated as planning artifacts only. A human study therefore remains uncollected.

\textbf{Locations:} Methods, p.~4, lines 173--182; Results and Fig.~3, pp.~8--9, lines 280--303; Discussion, p.~12, lines 411--416; Supplementary Note S5, pp.~8--9, lines 89--104.

\subsection*{4. Deterministic token-filter ablation}
\begin{reviewcomment}
The impact of the deterministic token filter is not isolated; compare generation with and without the filter.
\end{reviewcomment}

\Status{Substantially addressed}

A paired 48-payload, three-model no-filter contrast was added. Mean log-probability difference was -0.001797 (95\% CI, -0.037738 to 0.035394), and effective-rate difference was -0.062969 bits per full token (CI, -0.134265 to 0.007651); both selected intervals include zero. The Mistral roundtrip-stable-filter cell was unavailable for 48 work units because its isolated eligible vocabulary was empty. That cell is reported as unavailable, not failed.

\textbf{Locations:} Methods, pp.~2--3, lines 100--106; Results and Table~2, p.~9, lines 315--322; Supplementary Note S1, p.~3, lines 38--43; Supplementary Note S6 and Table~S10, pp.~8--11, lines 105--125.

\subsection*{5. Lead-in lengths and boundary failures}
\begin{reviewcomment}
The failed lead-in variant suggests boundary vulnerabilities; test several lead-in lengths and analyze the failure.
\end{reviewcomment}

\Status{Substantially addressed}

Lead-in lengths 2, 4, 8, 16, and 32 were evaluated against no lead-in. The compact 32-token contrast was adverse: mean log probability changed by -1.652327 and full-message length by +159.590278 tokens. The result is labeled exploratory/post-outcome and interpreted as a boundary/context association rather than a proven single mechanism. Greedy lead-in replay itself recovered 144/144 in the frozen robustness test, which distinguishes reproducible lead-ins from visible-text boundary stability.

\textbf{Locations:} Results and Table~2, p.~9, lines 305--314; Supplementary Note S1, p.~2, lines 29--36; Supplementary Note S6 and Table~S10, pp.~8--11, lines 105--125.

\subsection*{6. Dynamic semantic stopping}
\begin{reviewcomment}
Fixed tail caps may produce unnatural endings; implement and evaluate a dynamic semantic-completeness stopping rule.
\end{reviewcomment}

\Status{Substantially addressed}

The canonical segmented protocol now uses a frozen dynamic stopping rule requiring minimum tail length, declared terminal punctuation or double newline, balanced delimiters, and absence of specified dangling forms, with an emergency cap. It was compared with fixed 20--60-token sentence tails and no-tail generation. No-tail rate gains are described as denominator effects, not quality improvements. Because no participant study was conducted, the rule is explicitly called a heuristic rather than evidence of semantic completeness to readers.

\textbf{Locations:} Methods, p.~2, lines 87--99; Results, p.~9, lines 323--331; Supplementary Note S1, p.~2, lines 29--36; Supplementary Note S6, p.~8, lines 122--125.

\subsection*{7. Stronger neural steganalysis}
\begin{reviewcomment}
The preliminary detector is inadequate; add strong neural steganalysis models.
\end{reviewcomment}

\Status{Fully addressed}

The revision adds a TextCNN-equivalent linguistic-steganalysis architecture and a pretrained DeBERTa-v3-base classifier over a balanced 15,840-row corpus. Each architecture completed one matched, 18 held-out-template, three leave-one-model, and six leave-one-codec fits (56/56 total; 55,440 prediction rows). Matched ROC-AUC was 0.990190 and 0.999834, respectively. These near-perfect results are presented prominently as adverse evidence. Exact identity leakage checks passed; a 53-pair lexical near-duplicate limitation, exclusion sensitivity, and same-device CUDA repeatability are disclosed. No CPU/GPU equivalence is claimed.

\textbf{Locations:} Methods, pp.~4--5, lines 184--199; Results and Fig.~4, pp.~10--11, lines 340--370; Discussion, pp.~11--12, lines 405--410; Supplementary Note S7 and Fig.~S3, pp.~12--15, lines 134--159.

\subsection*{8. Single-topic versus multi-topic inference}
\begin{reviewcomment}
The topic difference lacks statistical validation; test it across broader prompt categories.
\end{reviewcomment}

\Status{Substantially addressed}

The primary prompt schedule now covers 18 templates in six categories. Fifteen locked within-model topic contrasts were extracted with Holm adjustment. One Llama artifact-count ratio excluded the null (1.276361; 95\% CI, 1.114743--1.461409; adjusted \(p=0.003294\)); the other 14 adjusted values were at least 0.396248. Because the compact extraction occurred after outcomes were available, it is labeled exploratory/post-outcome and does not support a general topic advantage.

\textbf{Locations:} Methods, p.~4, lines 146--160 and p.~5, lines 200--213; Results, p.~9, lines 332--339; Supplementary Note S6 and Table~S11, p.~12, lines 128--132.

\subsection*{9. Capacity--quality formalization}
\begin{reviewcomment}
The capacity--quality relationship is visualized but not mathematically formalized; define expected bounds between payload capacity and cover length.
\end{reviewcomment}

\Status{Substantially addressed}

The revision defines codec-implied forced length, nominal and effective rate, segmentation/tail length identities, a bounded-rank proposition, and a conditional forced-token probability bound. Empirical checking covered 17,424 records, with 7,008 fully evaluable and zero forced-position, rate-bound, or supported quality-bound violations; 2,976 positive residuals quantify tail/cover overhead. The complete encoder--decoder cascade is not directly trace-evaluable from the retained record schema, and model likelihood is not presented as proof of human quality.

\textbf{Locations:} Methods, p.~3, lines 124--135; Results and Fig.~1, pp.~6--7, lines 239--251; Supplementary Note S8 and Table~S14, pp.~14--15, lines 160--194.

\subsection*{10. Cryptographic payload realism}
\begin{reviewcomment}
Ciphertext-like base64 strings were not generated by real encryption algorithms; use standard cryptographic protocols.
\end{reviewcomment}

\Status{Fully addressed within a public test-vector design}

The corpus now contains deterministic outputs from declared algorithms: SHA-256, HMAC-SHA256, 96-bit nonces, 128-bit tokens, UUIDv4, AES-256-GCM, ChaCha20--Poly1305, and Ed25519. The encryption and signature outputs use their declared base64 serializations. They are reproducible public test vectors, not operational secrets or deployment payloads; that boundary is explicit.

\textbf{Locations:} Methods, p.~4, lines 138--145; Discussion, p.~12, lines 417--422; Supplementary Note S2 and Table~S1, p.~3, lines 54--60; worked example, p.~18, lines 205--220.

\subsection*{11. Multiple language-model families}
\begin{reviewcomment}
Dependence on one architecture limits generalizability; replicate with at least one additional open-source model.
\end{reviewcomment}

\Status{Substantially addressed}

The completed primary study uses three open-source 7B--8B families: Meta-Llama-3-8B-Instruct, Qwen2.5-7B-Instruct, and Mistral-7B-Instruct-v0.3, all under pinned Q4\_K\_M configurations. Recovery and principal summaries are stratified by model. The paper does not extend this evidence to very-large proprietary models or changing hosted services; the configured GGUF files are also absent from the present checkout, so fresh local generation replay is unavailable.

\textbf{Locations:} Methods, p.~4, lines 146--160; Results, pp.~5--6, lines 220--236; Discussion, p.~12, lines 417--422; Supplementary Notes S2--S3, pp.~3--6, lines 58--78.

\subsection*{12. Computational overhead}
\begin{reviewcomment}
Add time and memory analysis for RankCloak encoding and decoding.
\end{reviewcomment}

\Status{Substantially addressed}

Across 18 primary model--protocol cells and 12,144 matched runtime rows, mean inclusive generation time ranged 1.674368--16.087628 seconds, payload throughput 13.631855--100.851772 bits/s, and effective rate 0.710004--7.322003 bits per full token. The Supplement reports the complete plot/table, component scopes, detector CUDA time, and peak allocated VRAM. Measurements are inclusive wrapper wall time; CPU time and repeated warm-up measurements were unavailable, so no broad hardware or CPU/GPU claim is made.

\textbf{Locations:} Methods, p.~5, lines 214--217; Results and Table~2, pp.~11--12, lines 373--390; Supplementary Note S9, Fig.~S4, and Tables~S15--S16, pp.~15--18, lines 195--204.

\subsection*{13. Reliance on preprints and technical reports}
\begin{reviewcomment}
The reference list relies heavily on unrefereed preprints, including original references 1, 12, 28, 32, and 34; replace them with peer-reviewed versions where available or reduce reliance.
\end{reviewcomment}

\Status{Fully addressed where a verified archival version exists}

Every cited entry was audited. The former Calgacus and \emph{Early Signs} preprints now cite their ICLR 2026 proceedings versions; DAIRstega now cites the 2025 \emph{Applied Soft Computing} article; and the AI-generated-text detection citation now uses its 2025 \emph{Transactions on Machine Learning Research} version. Tangential preprints that were not needed for the revised argument were removed. Standards/RFCs remain where they are the authoritative algorithm specifications. No bibliographic entry was invented; verification evidence and the disposition of original references 1, 12, 28, 32, and 34 are recorded in \path{paperV2/CITATION_AUDIT.md}.

\textbf{Locations:} Related work, pp.~1--2, lines 1--55; References, pp.~13--15, lines 470--553; citation audit file.

\section*{Reviewer 2}

\subsection*{Positive assessment}
\begin{reviewcomment}
The reviewer recognized the separation of payload representation, recovery correctness, and cover quality; the comparison of direct, bounded, and segmented variants; and the explicit statement of assumptions and failure cases.
\end{reviewcomment}

\Status{Retained and strengthened}

The core RankCloak formulation remains the manuscript's contribution. The revised work preserves direct subword, bounded ASCII-byte and hex-nibble representations, non-segmented and segmented generation, forced-span decoding, and non-payload tails, while replacing pilot evidence and sharpening the replay and claim boundaries.

\textbf{Locations:} Introduction, pp.~1--2, lines 1--55; Methods, pp.~2--5, lines 58--217; Supplementary Note S1, pp.~1--3, lines 5--52.

\subsection*{1. Scale, models, languages, and deployment}
\begin{reviewcomment}
The study is controlled and limited; investigate larger-scale LLMs, different model families, multilingual settings, realistic payloads, and real-world deployment.
\end{reviewcomment}

\Status{Partially addressed}

Scale, family diversity, payload realism, and secondary-language exact-copy recovery were expanded: 480 algorithm-defined artifacts, three model families, 18 English templates, and 288 Spanish plus 288 Simplified Chinese trials. All 576 secondary-language trials recovered under saved-token replay. These tests do not establish multilingual naturalness. Very-large proprietary models and real-world deployment were not tested, and no deployment claim is made.

\textbf{Locations:} Methods, p.~4, lines 138--160; Results and Table~1, pp.~5--6, lines 220--236; computational/generalizability Results, pp.~11--12, lines 373--390; Discussion and limitations, pp.~11--12, lines 392--461; Supplementary Notes S2--S3, pp.~3--6, lines 54--78.

\subsection*{2. Human perception, steganalysis, and suggested Ding et al. citation}
\begin{reviewcomment}
Model-based metrics alone are insufficient; add human evaluation and stronger steganalysis, and compare with linguistic-steganography work. The reviewer suggests Ding et al., \emph{Self-Attention with Cross-Lingual Position Representation}, if its representation-learning connection can be clarified.
\end{reviewcomment}

\Status{Partially addressed for human evaluation; fully addressed for neural detection; citation clarified}

Automated readability and held-out model diagnostics were added, but no participant study was performed; the revision says so explicitly. Neural steganalysis was expanded to TextCNN-equivalent and DeBERTa-v3-base architectures over matched and three held-out regimes, with adverse near-perfect matched AUC, leakage audits, and sensitivity analysis. The literature review now includes directly relevant linguistic steganography, linguistic steganalysis, and generated-text detection work.

We verified the suggested Ding et al. paper. It concerns cross-lingual position representations for machine translation, not a linguistic-steganography method or detector. Citing it as a baseline would imply a relationship that the present experiments do not test. We therefore did not force the citation and instead clarified RankCloak's relationship to directly relevant rank transfer and token-selection literature.

\textbf{Locations:} Introduction, pp.~1--2, lines 1--55; Methods, pp.~4--5, lines 173--199; automated-quality Results, pp.~8--9, lines 280--303; neural-detection Results, pp.~10--11, lines 340--370; Supplementary Notes S5 and S7, pp.~8--9 and 12--15, lines 89--104 and 134--159; citation audit file.

\section*{Summary of remaining boundaries}

The following requests are not represented as completed experiments: participant-based readability evaluation, real-world deployment, very-large proprietary-model evaluation, broad multilingual naturalness, and the untested perturbation classes. CPU timing, repeated warm-up measurements, fresh generation replay without the configured GGUF files, and a directly trace-evaluable complete encoder--decoder cascade also remain unavailable. DOI-bearing archival deposition is an author action if required at submission. These boundaries are stated in the manuscript rather than hidden in this response.

\end{document}
